\documentclass[11pt,a4paper]{article}

\usepackage[margin=2.5cm]{geometry}
\usepackage{graphicx}
\usepackage{amsmath}
\usepackage{amssymb}
\usepackage{booktabs}
\usepackage{microtype}
\usepackage{cite}
\title{\textbf{Dual-Morphology Spinodal Demixing in Non-Classical Calcium Carbonate Precipitation: Coexistence of an Extended Viscoelastic Fractal Mesh and Capillary-Pinched Dense Liquid Droplets}}

\author{Tomasz M. Stawski\thanks{Email: \texttt{tomasz.stawski@gmail.com}} \\
\small Federal Institute for Materials Research and Testing (BAM), Richard-Willst{\"a}tter-Str. 11, 12489 Berlin, Germany}

\date{\today}

\begin{document}

\maketitle

\begin{abstract}
During the rapid precipitation of calcium carbonate at high supersaturations and alkaline pH, the system bypasses classical monomer-by-monomer nucleation. Crucially, even in regimes where solid amorphous calcium carbonate (ACC) nanoparticles or crystalline polymorphs (vaterite, calcite) do not formally precipitate, in situ time-resolved Small-Angle and Wide-Angle X-ray Scattering (SAXS/WAXS) reveals the concurrent emergence of an extended, space-filling continuous fractal mesh ($D_m \approx 1.83 - 1.95$) and a prominent ultra-low-$q$ scattering signature corresponding to giant spherical microstructures ($\langle R \rangle \approx 3.2\ \mu\mathrm{m}$). Here, we provide an analytical scattering model and physical framework resolving this coexistence. We demonstrate that the ``large spheres'' are not solid crystals or rigid ACC colloids, but rather surface-tension-rounded, dense liquid-like precursor droplets (emulsion domains) produced via Rayleigh--Plateau capillary pinch-off of cylindrical spinodal struts and Viscoelastic Phase Separation (VPS). Supported by in situ WAXS confirming the total absence of Bragg diffraction, and taking into account the $V^2 \propto R^6$ forward-scattering weighting, our model proves that a minor volume fraction ($\sim 8\%$) of capillary-detached droplets naturally dominates the ultra-low-$q$ window while leaving the bulk continuous gel network ($\sim 92\%$ by volume) to govern intermediate-to-high-$q$ fractal scattering.
\end{abstract}

\section{Introduction and Physical Premise}

Non-classical pathways in mineral formation have transformed our understanding of crystallization in synthetic, geological, and biological systems \cite{deyoreo2015crystallization, gebauer2008stable, stawski2016formation, stawski2019nucleation}. In highly supersaturated aqueous solutions containing $\mathrm{Ca}^{2+}$ and $\mathrm{CO}_3^{2-}$ ions, precipitation is frequently initiated through liquid--liquid phase separation (LLPS) \cite{gower2000biomimetic, gower2008biomimetic, wallace2013microscopic, bewernitz2012metastable}. Deep within the spinodal instability boundary, barrierless demixing drives the spontaneous condensation of hydrated ionic polymers or Dynamically Ordered Liquid-Like Oxyanion Polymers (DOLLOPs) \cite{cahn1965phase, demichelis2011stable}. 

In many experimental regimes, this spinodal demixing leads to an extended, percolating continuous network (a bicontinuous gel-like mesh) without the formal precipitation of crystalline phases (calcite, aragonite, vaterite) or solid amorphous calcium carbonate (ACC) nanoparticles. Yet, in situ Small-Angle X-ray Scattering (SAXS) data collected over broad reciprocal-space windows ($q \in [0.002, 3.5]\ \mathrm{nm}^{-1}$) systematically exhibit two ostensibly conflicting structural motifs:
\begin{enumerate}
    \item \textbf{At intermediate to high scattering vectors ($0.005 \le q \le 0.2\ \mathrm{nm}^{-1}$):} An extended, scale-invariant mass fractal power-law decay ($I(q) \propto q^{-D_m}$, with $D_m \approx 1.83 - 1.95$) reflecting a tenuous, space-filling continuous mesh formed by Diffusion-Limited Cluster Aggregation (DLCA) of primary hydrated nodes ($r_0 \approx 8 - 9\ \mathrm{nm}$).
    \item \textbf{At ultra-low scattering vectors ($q < 0.005\ \mathrm{nm}^{-1}$):} A steep, asymptotic intensity surge obeying Porod's law ($I(q) \propto q^{-4}$) accompanied by deep, undamped Bessel interference fringes characteristic of giant, highly monodisperse spherical domains with mean radii $\langle R \rangle \approx 2.9 - 3.5\ \mu\mathrm{m}$.
\end{enumerate}

The physical origin of these ``large spheres'' within a non-crystallizing, spinodal continuous-mesh reaction has remained an active puzzle. Here, we demonstrate that their presence is neither anomalous nor indicative of premature crystallization. Instead, it is the direct, hydrodynamic consequence of capillary instabilities (Rayleigh--Plateau break-up) and Viscoelastic Phase Separation (VPS) in a bicontinuous, dynamic liquid--liquid mixture.

\section{Multi-Scale Scattering Model}

The total coherent scattering intensity $I_{\mathrm{tot}}(q)$ across the full experimental range (spanning camera lengths of $2.0\ \mathrm{m}$ and $31.0\ \mathrm{m}$ at ESRF ID02 \cite{narayanan2018id02}) is formulated as an incoherent superposition of three decoupled structural regimes plus a constant solvent background:
\begin{equation}
I_{\mathrm{tot}}(q) = I_{\mathrm{drop}}(q) + I_{\mathrm{mesh}}(q) + I_{\mathrm{fluc}}(q) + I_{\mathrm{bkg}}
\label{eq:Itot}
\end{equation}

\subsection{Tier 1: Giant Dense-Liquid Droplets ($q < 0.005\ \mathrm{nm}^{-1}$)}

The low-$q$ domain is governed by isolated, surface-tension-rounded spherical droplets of dense liquid precursor. Assuming a continuous lognormal size distribution $f_N(R)$ defined by median radius $R_0$ and geometric standard deviation $\sigma$, the scattering intensity is:
\begin{equation}
I_{\mathrm{drop}}(q) = n_{\mathrm{drop}} (\Delta \rho_{\mathrm{drop}})^2 \int_0^\infty V(R)^2 P_{\mathrm{sph}}(q, R) f_N(R) \, dR
\label{eq:Idrop}
\end{equation}
where $V(R) = \frac{4}{3}\pi R^3$, $\Delta \rho_{\mathrm{drop}} = \rho_{\mathrm{drop}} - \rho_{\mathrm{water}}$ is the electron scattering length density (SLD) contrast, and $P_{\mathrm{sph}}(q, R)$ is the isotropic sphere form factor:
\begin{equation}
P_{\mathrm{sph}}(q, R) = \left[ 3 \frac{\sin(qR) - qR\cos(qR)}{(qR)^3} \right]^2 = \left[ \frac{3 j_1(qR)}{qR} \right]^2
\end{equation}

In terms of the macroscopic droplet volume fraction $\phi_{\mathrm{drop}} = n_{\mathrm{drop}} \langle V \rangle$:
\begin{equation}
I_{\mathrm{drop}}(q) = \phi_{\mathrm{drop}} (\Delta \rho_{\mathrm{drop}})^2 \frac{\langle V^2 P_{\mathrm{sph}}(q) \rangle}{\langle V \rangle}
\end{equation}

In the Porod asymptotic regime ($q R \gg 1$), destructive interference dampens the Bessel oscillations, yielding Porod's universal law for 3D compact objects with sharp interfaces \cite{porod1951die}:
\begin{equation}
\lim_{qR \gg 1} I_{\mathrm{drop}}(q) = 2\pi (\Delta \rho_{\mathrm{drop}})^2 \frac{S_{\mathrm{tot}}}{q^4}
\label{eq:porod}
\end{equation}
where $S_{\mathrm{tot}} = 4\pi \langle R^2 \rangle n_{\mathrm{drop}}$ is the total specific interfacial area. Crucially, the observed power-law exponent of $\partial \ln I / \partial \ln q \approx -4.0$ to $-5.0$ at $q \in [0.002, 0.004]\ \mathrm{nm}^{-1}$ mathematically requires a smooth, 3D compact interface ($D=3$), which cannot be produced by any single fractal structure.

\subsection{Tier 2: Continuous Spinodal Mesh ($0.005 \le q \le 0.2\ \mathrm{nm}^{-1}$)}

The bulk of the mineral mass resides in the continuous, interconnected spinodal network. This structure is modeled as an extended mass fractal network composed of primary hydrated nodes of mean radius $r_0 \approx 8.8\ \mathrm{nm}$ interlinked with fractal dimension $D_m$ and upper correlation length $\xi$ \cite{teixeira1988small, sorensen2001light}:
\begin{equation}
I_{\mathrm{mesh}}(q) = \phi_{\mathrm{mesh}} (\Delta \rho_{\mathrm{mesh}})^2 V_{\mathrm{node}} P_{\mathrm{node}}(q, r_0) S_{\mathrm{frac}}(q, D_m, \xi)
\label{eq:Imesh}
\end{equation}
where $P_{\mathrm{node}}(q, r_0)$ is the polydisperse spherical form factor of the nodes, and $S_{\mathrm{frac}}(q)$ is the Teixeira--Sorensen structure factor:
\begin{equation}
S_{\mathrm{frac}}(q) = 1 + \frac{D_m \Gamma(D_m - 1)}{(q r_0)^{D_m} [1 + (q\xi)^{-2}]^{\frac{D_m - 1}{2}}} \sin\left[ (D_m - 1) \arctan(q\xi) \right]
\label{eq:Sfrac}
\end{equation}

In the fractal power-law regime ($1/\xi \ll q \ll 1/r_0$), Eq.~(\ref{eq:Sfrac}) reduces directly to:
\begin{equation}
I_{\mathrm{mesh}}(q) \propto q^{-D_m} \approx q^{-1.84}
\end{equation}
which is the signature of 3D Diffusion-Limited Cluster Aggregation (DLCA) and percolation of spinodal channels.

\subsection{Tier 3: Internal Mesh Fluctuations and Prenucleation Clusters ($q > 0.2\ \mathrm{nm}^{-1}$)}

At high $q$, the intensity levels off into an elevated plateau above the flat incoherent solvent background. This is described via two complementary physical pictures:
\begin{enumerate}
    \item \textbf{Prenucleation Clusters (DOLLOPs):} A population of sub-nanometer hydrated ion associates ($\langle R_{\mathrm{PNC}} \rangle = 0.53 \pm 0.15\ \mathrm{nm}$, $\phi_{\mathrm{PNC}} \approx 1.84 \times 10^{-3}$) residing within the interstitial water channels \cite{gebauer2008stable, demichelis2011stable}.
    \item \textbf{Viscoelastic Network Fluctuations:} Dynamic concentration fluctuations of the polymer-like network with correlation length $\xi_{\mathrm{fluc}}$ and Porod surface exponent $\mu = 4.0$ according to the Fisher--Burford / Tanaka framework \cite{tanaka2000viscoelastic, fisher1967theory}:
    \begin{equation}
    I_{\mathrm{fluc}}(q) = \frac{A_{\mathrm{fluc}}}{\left[ 1 + \frac{2}{3\mu} (q \xi_{\mathrm{fluc}})^2 \right]^{\mu / 2}} = \frac{A_{\mathrm{fluc}}}{\left[ 1 + \frac{1}{6} (q \xi_{\mathrm{fluc}})^2 \right]^2}
    \end{equation}
    where $\xi_{\mathrm{fluc}} = 4.62 \pm 0.64\ \mathrm{nm} \approx \frac{1}{2} r_0$, representing the transient mesh size between physical crosslinks of the continuous gel.
\end{enumerate}

\section{Physical Origin of Large Spheres}

\subsection{Rayleigh--Plateau Instability of Bicontinuous Struts}

In classical Cahn--Hilliard spinodal decomposition, liquid--liquid demixing initiates through long-wavelength concentration fluctuations, forming an isotropic bicontinuous network of dense liquid channels \cite{cahn1965phase}. Because the dense liquid phase possesses a finite interfacial tension $\gamma \sim 1 - 10\ \mathrm{mN/m}$ against water, any cylindrical ligament of radius $R_{\mathrm{strut}}$ is unstable to axisymmetric capillary undulations:
\begin{equation}
r(z, t) = R_{\mathrm{strut}} + \epsilon_0 \cos(k z) e^{\omega t}
\end{equation}

According to Rayleigh and Plateau \cite{plateau1873statique, rayleigh1878instability}, the surface area per unit length decreases if the perturbation wavelength $\lambda = 2\pi / k$ exceeds the circumference of the cylinder:
\begin{equation}
\lambda > 2\pi R_{\mathrm{strut}} \quad \iff \quad k R_{\mathrm{strut}} < 1
\end{equation}

For a viscous fluid thread embedded in an external viscous medium, the hydrodynamic dispersion relation $\omega(k)$ exhibits a fastest-growing unstable mode at:
\begin{equation}
k_{\mathrm{opt}} R_{\mathrm{strut}} \approx 0.697 \implies \lambda_{\mathrm{opt}} \approx 9.02 R_{\mathrm{strut}}
\end{equation}

Assuming volume conservation during capillary collapse, one segment of length $\lambda_{\mathrm{opt}}$ contracts into a single spherical droplet of radius $R_{\mathrm{drop}}$:
\begin{equation}
\frac{4}{3}\pi R_{\mathrm{drop}}^3 = \pi R_{\mathrm{strut}}^2 \lambda_{\mathrm{opt}} \implies R_{\mathrm{drop}} \approx 1.89 R_{\mathrm{strut}}
\label{eq:Rdrop}
\end{equation}

Local shear gradients during rapid reactant mixing yield spinodal ligament radii up to $R_{\mathrm{strut}} \approx 1.5 - 1.7\ \mu\mathrm{m}$. Eq.~(\ref{eq:Rdrop}) predicts pinched-off spherical droplet radii of:
\begin{equation}
R_{\mathrm{drop}} \approx 2.8 - 3.2\ \mu\mathrm{m}
\end{equation}
in quantitative agreement with the experimentally fitted value ($\langle R_{\mathrm{drop}} \rangle = 3.21 \pm 0.04\ \mu\mathrm{m}$).

\subsection{Viscoelastic Phase Separation and Droplet Coexistence}

Because the dense mineral network is composed of transiently crosslinked ion-polymer chains, it exhibits pronounced \emph{dynamic asymmetry}: the structural relaxation time of the network $\tau_{\mathrm{mesh}}$ is orders of magnitude slower than that of the solvent $\tau_{\mathrm{solv}}$ ($\tau_{\mathrm{mesh}} \gg \tau_{\mathrm{solv}}$). Under Tanaka's Viscoelastic Phase Separation (VPS) framework \cite{tanaka2000viscoelastic, tanaka2001viscoelastic}:
\begin{enumerate}
    \item \textbf{Transient Elastic Stress:} During coarsening, capillary forces attempt to shrink the interconnected network, inducing mechanical tension $\sigma_{ij}$ along the struts.
    \item \textbf{Stress-Induced Rupture:} When the accumulated elastic energy exceeds the transient yield threshold of the fluid ligaments, the struts rupture mechanically. The detached segments rapidly contract under surface tension into compact spheres.
    \item \textbf{Dynamic Arrest:} Simultaneously, network percolation in low-stress regions arrests into an extended, space-filling fractal mesh ($D_m \approx 1.84 - 1.93$).
\end{enumerate}

Thus, the system naturally divides into a dual morphology: an arrested continuous mass fractal gel coexisting with an emulsion of capillary-pinched dense liquid globules.

\subsection{Absence of Bragg Peaks in in situ WAXS}

To confirm whether the micrometer spheres contain crystalline vaterite or calcite, in situ Wide-Angle X-ray Scattering (WAXS) was acquired simultaneously with SAXS. 

Throughout the entire observational window ($t = 0 - 7.0\ \mathrm{s}$), the WAXS diffractograms exhibit no sharp Bragg reflections corresponding to vaterite [e.g., (110) at $d = 3.58\ \text{\AA}$, (112) at $d = 3.29\ \text{\AA}$, (114) at $d = 2.73\ \text{\AA}$] or calcite [e.g., (104) at $d = 3.03\ \text{\AA}$]. Instead, the pattern displays only the broad, diffuse amorphous liquid-structure halo of bulk water and hydrated Ca-carbonate ion associations centered near $q \approx 20\ \mathrm{nm}^{-1}$. 

This directly confirms that the large spheres are purely liquid-like precursor droplets (PILP domains) rather than crystalline spherulites.

\section{The Forward Scattering Disparity ($V^2 \propto R^6$)}

A central experimental question is: \emph{If the continuous mesh forms over $90\%$ of the precipitated mass, why do the large spheres dominate the ultra-low-$q$ SAXS profile?}

In small-angle scattering, the zero-angle forward scattering intensity $I(0)$ of an individual scattering object scales with the square of its volume:
\begin{equation}
I_{\mathrm{single}}(0) = (\Delta \rho)^2 V^2 \propto R^6
\end{equation}

Comparing a single micrometer droplet ($R = 3210\ \mathrm{nm}$) to a primary network node ($r_0 = 8.8\ \mathrm{nm}$):
\begin{equation}
\frac{V_{\mathrm{drop}}}{V_{\mathrm{node}}} = \left(\frac{3210}{8.8}\right)^3 \approx 4.85 \times 10^7
\end{equation}
\begin{equation}
\frac{I_{\mathrm{drop}}(0)}{I_{\mathrm{node}}(0)} = \left(\frac{V_{\mathrm{drop}}}{V_{\mathrm{node}}}\right)^2 \approx \mathbf{2.35 \times 10^{15}}
\end{equation}

A single $3.2\ \mu\mathrm{m}$ droplet scatters with the forward intensity of more than two quadrillion primary network nodes!

Macroscopically, for droplet volume fraction $\phi_{\mathrm{drop}} \approx 7.0 \times 10^{-4}$ ($\approx 8.2\%$ of total mineral mass) and mesh volume fraction $\phi_{\mathrm{mesh}} \approx 7.8 \times 10^{-3}$ ($\approx 91.8\%$ of total mineral mass):
\begin{equation}
\frac{I_{\mathrm{drop}}(0)}{I_{\mathrm{mesh}}(0)} = \frac{\phi_{\mathrm{drop}} V_{\mathrm{drop}}}{\phi_{\mathrm{mesh}} V_{\mathrm{cluster}}} \approx \mathbf{58.9}
\end{equation}

Consequently, at $q_{\mathrm{min}} = 0.00213\ \mathrm{nm}^{-1}$, the capillary droplets account for \textbf{$77.8\%$} of the total measured intensity. However, because $I_{\mathrm{drop}}(q)$ drops precipitously as $q^{-4}$, its contribution falls below $2\%$ by $q \approx 0.008\ \mathrm{nm}^{-1}$. At this point, the scattering of the continuous mass fractal mesh ($I \propto q^{-1.84}$) emerges without distortion across the intermediate $q$ range.

\section{Experimental Validation}

Table~\ref{tab:params} presents the quantitative parameter set obtained by non-linear least-squares fitting of Eq.~(\ref{eq:Itot}) to the calibrated, merged dataset of sample ND-i-009 ($40.0\ \mathrm{mM}\ \mathrm{CaCO}_3$, $\mathrm{pH}\ 10.0$).

\begin{table}[htbp]
\centering
\caption{\label{tab:params}Calibrated multi-scale model parameters for sample ND-i-009 ($40.0\ \mathrm{mM}$, $\mathrm{pH}\ 10.0$) across merged short and long camera distances.}
\vspace{2mm}
\begin{tabular}{llcc}
\toprule
\textbf{Component} & \textbf{Parameter} & \textbf{Fitted Value} & \textbf{Std. Error} \\
\midrule
Capillary Droplets & $\langle R_{\mathrm{drop}} \rangle$ & $3.21\ \mu\mathrm{m}$ & $\pm 0.04\ \mu\mathrm{m}$ \\
& $\sigma_{\mathrm{drop}}$ & $0.082$ & $\pm 0.006$ \\
& $\phi_{\mathrm{drop}}$ & $7.02 \times 10^{-4}$ & $\pm 0.15 \times 10^{-4}$ \\
\midrule
Continuous Mesh & $D_m$ & $1.925$ & $\pm 0.002$ \\
& $\langle r_0 \rangle$ & $9.20\ \mathrm{nm}$ & $\pm 0.07\ \mathrm{nm}$ \\
& $\sigma_0$ & $0.312$ & $\pm 0.008$ \\
& $\phi_{\mathrm{mesh}}$ & $7.81 \times 10^{-3}$ & $\pm 0.12 \times 10^{-3}$ \\
\midrule
Fluctuations/PNCs & $\langle R_{\mathrm{PNC}} \rangle$ & $0.528\ \mathrm{nm}$ & $\pm 0.147\ \mathrm{nm}$ \\
& $\phi_{\mathrm{PNC}}$ & $1.84 \times 10^{-3}$ & $\pm 0.22 \times 10^{-3}$ \\
& $\xi_{\mathrm{fluc}}$ & $4.62\ \mathrm{nm}$ & $\pm 0.64\ \mathrm{nm}$ \\
\midrule
Background & $bkg$ & $0.0216\ \mathrm{cm}^{-1}$ & $\pm 0.0003\ \mathrm{cm}^{-1}$ \\
\midrule
Goodness of Fit & $\chi^2_{\mathrm{red}}$ (Overall) & $\mathbf{1.958}$ & --- \\
& $\chi^2_{\mathrm{red}}$ ($q > 0.3$) & $\mathbf{0.542}$ & --- \\
& AIC & $\mathbf{1202.8}$ & --- \\
\bottomrule
\end{tabular}
\end{table}

\section{Conclusions}

The observation of a pronounced, oscillating spherical form factor at ultra-low $q$ in the absence of crystalline polymorphs or solid ACC nanoparticles is completely explained by the fluid mechanics and thermodynamics of spinodal demixing:
\begin{enumerate}
    \item The ``large spheres'' are \textbf{dense liquid precursor droplets} formed via the \textbf{Rayleigh--Plateau capillary instability} of cylindrical bicontinuous network struts and stress-relaxation during \textbf{Viscoelastic Phase Separation}.
    \item Because these domains are liquid, interfacial tension forces them into smooth spheres with molecularly sharp boundaries, generating the observed $q^{-4}$ Porod slope and sharp Bessel interference fringes.
    \item While the continuous fractal mesh accounts for over \textbf{$90\%$} of the total precipitated mineral mass, the $V^2 \propto R^6$ forward-scattering weighting ensures that the remaining \textbf{$\sim 8\%$} droplet population naturally dominates the ultra-low-$q$ window.
\end{enumerate}

\bibliographystyle{unsrt}
\bibliography{references}

\end{document}
